\part{Unchecked identifications and explicit limits of the result}
\section{The remaining gate is an identification theorem, not the trace formula}
\paragraph{U-STEP-14 (historical source matching).}
The note's step 1 names a Stepanov determinant, and step 4 names $W(1_S)\ge0$ via a determinant; neither gives a matrix. The full-source exact-name and independent structural searches do not recover a unique declaration for either phrase. Archive metadata and the public mirror history supply no October chronology. The exact candidate Weil-form calculation in \fid{STEP6-01--03} reconstructs the two-indicator test, its scalar prime-sum cross term and its half-plane consequence; it narrows the remaining dictionary to normalization, linear versus quadratic convention, the sign/amplitude and the actual character/window specialization. The step-1 search read actual modular Wronskians, theta evaluation determinants, finite monomial interpolation matrices, directional Hasse derivatives of polynomial minors, and Frobenius/Moore determinants. The step-4 search distinguished finite Pl\"ucker indicator determinants, alternating Weil-named forms, quotient measures, positive Rankin--Selberg integrands and the finite inertia theorem of \fid{DET-10}. Exact locators and read scope are preserved in the evidence notes. Each inspected candidate has a specified finite algebraic or analytic domain; resemblance does not establish the missing historical correspondence. Step 7 is now supported by actual finite minor, exterior-power and period identities, but the note does not identify every determinant with one declaration. This remains an exact-source correspondence gap. The audit-constructed interlacing selector and the actual rank-one period factorization in \fid{STEP7-01--02} further distinguish the finite mechanisms. A full formal proof is not rechecked or rejected on that account.

\paragraph{U-LIMIT (the stronger possible interpretation).}
The literal Fable proposal and the fixed finite polynomial identification are refuted. It remains unproved that a transformed family of Alp\"oge's determinants is a compatible family of compressions or approximants of \fid{B-01}. No impossibility theorem about every such future transformation is asserted. The current sufficient gate is concrete: exhibit the actual operators on a common Hilbert space; establish their Hermitian positivity and trace-norm convergence; and prove their limiting spectral determinant equals $G$ independently of assuming RH. The checked \fid{DET-05} then finishes the desired implication. Alternatively provide a positive arithmetic kernel and a unitary/spectral identification theorem.

The approaches tried were: the supplied note and map; exact entry-file reading; a literal finite Fourier identification (refuted by eigenvalues); path-to-trace identification (refuted and replaced by the rank-one variation); fixed-degree determinant identification (refuted by infinitely many zeros); actual Gram/Capelli/exterior-period source analysis; and the scalar zero-free endpoint trace. A global textual search for \texttt{Fredholm}, \texttt{riemannXi} and \texttt{xiFunction} returns no occurrences, while trace-class automorphic operators do occur. This is supporting search evidence only, not a proof that an unnamed equivalent construction is impossible.

\paragraph{U-MAS (conditional source claims remain conditional).}
The proof of the complete Masullo draft is not a finding of this run. The actual single-operator Fredholm bridge is now proved in \fid{MAS-10}, conditional on its full spectral-jet quotient claim; there is no remaining trace-class gap in that construction. The draft's stated trace-class/arithmetic formulas and general-Hecke GRH consequence retain their conditional labels, the $I$-quantifier qualification, the distinct Chapter 13 statement, and the boundary correction. We do not fill the empty Theorem 13.3.1.2 with an invented assertion. Determining the truth of those deep draft claims is a separate audit from identifying their statements and comparing their operator types.

\paragraph{U-TRANSCRIPTION (images needed).}
The mathematics specifies the corrections; attribution of each difference to the transcription rather than to Grothendieck's draft awaits the original images. The required next step is the second-pass comparison of the listed cote/scan/line entries, not another analytic proof.

All findings in Part I have complete checked proofs or are explicitly attributed conditional applications of cited results. Part II contains no new claimed theorem with an unacknowledged gap. The status file records the exact limits above rather than declaring a universal identification disproved by a failed text search. No three-hour or hourly movement review is claimed before its scheduled time; elapsed-time checkpoints are handled in the dated ledger.

\begin{thebibliography}{99}\raggedright
\bibitem{archive} A.~Grothendieck, unpublished archive, cote 1, supplied transcriptions: scans 78--80, formulas (15)--(23); scan 81, item 3; scan 110, lines 1--12; scan 111, lines 1--23; scan 112, lines 1--23; scan 113, lines 1--23. No manuscript image is part of the supplied ZIP.
\bibitem{groth56} A.~Grothendieck, \emph{La th\'eorie de Fredholm}, Bull. Soc. Math. France \textbf{84} (1956), 319--384. Chapter II, \S4, Theorem 2 (p.~354) and Theorem 4 (p.~359): invertibility and zero multiplicity. \url{https://www.numdam.org/item/10.24033/bsmf.1476.pdf}.
\bibitem{alpoge} L.~Alp\"oge, \emph{Remark}, supplied four-page note dated in the prompt 20 November 2026. Timeline \S2: steps 1, 3, 4 and footnote 3, p.~3; steps 7, 8, 16 and 20, p.~4. \url{https://alpo.ge/remark.pdf}. The theorem convention is that of the task; no Lean replay is claimed.
\bibitem{lean} Alp\"oge Lean sources: supplied \path{ALPOGE_Main.lean}, declarations \texttt{riemannHypothesis}, \texttt{riemannZeta\_eq\_zero\_re}, \texttt{LFunction\_ne\_zero\_of\_half\_lt\_re}, \texttt{LFunction\_eq\_zero\_re}, \texttt{LFunction\_eq\_zero\_re\_of\_isPrimitive}; propositions in \path{ALPOGE_Statements.lean}. Additional definition files and exact declarations are cited in the text and source manifest. Public mirror snapshot: \url{https://github.com/CocoSgt/riemann-hypothesis-lean/releases/tag/v2026-11-20}.
\bibitem{map} \emph{ALPOGE\_LEAN\_MAP\_20261120.md}, prior textual mapping run, Claude Opus 5.5, effort medium, 20 November 2026; \S3 dependency levels and \S4 step table. Textual mapping only; no build claimed by that source.
\bibitem{sa23} \emph{INTAKE\_RESEARCH\_SA23\_20261120.md}, supplied SA23 run, Claude Opus 5.5, effort max as identified in the task; \S1.H and \S4.1(1)--(4), especially \texttt{geo-H-bounds}, \texttt{geo-H-genus}, \texttt{geo-H-lidskii}, \texttt{geo-H-1956-reconstruction}. The present paper checks the stated $\Xi$ consequences and repairs their overbroad interpretations.
\bibitem{fable} \emph{FABLE\_PROGRAM\_SUGGESTIONS\_20261120.md} and \emph{FABLE\_PROGRAM\_SUGGESTIONS\_20261120\_EXPLAINED.md}, supplied Fable 5.2 run, effort max as identified in the task, suggestion 3 in each file. Toy theorem and proposed Fourier/Fredholm and Landau identifications.
\bibitem{masullo} A.~M.~Masullo, \emph{The Music of the Primes}, supplied 940-page draft, dated 20 November 2026 in the task. Exact statements used: Definition 10.1.1.1; Theorems 10.2.2.1, 10.3.1.1, 10.3.2.1, 12.5.1.1, 13.3.1.1; the empty heading 13.3.1.2; Definitions 10.4.1.1, 10.4.1.3, 10.4.3.7; Remark 10.4.1.4; norm formulas (9.2.2.34), (9.2.2.40)--(9.2.2.41); introductory \S\S1.2, 1.4.5--1.4.8, 1.5. All claimed general-Hecke consequences remain conditional on the draft.
\bibitem{titch} E.~C.~Titchmarsh, \emph{The Theory of the Riemann Zeta-Function}, 2nd ed., revised by D.~R.~Heath-Brown, Oxford, 1986. \S2.1, Theorem 2.12 and (2.12.4)--(2.12.5), pp.~29--30; Theorem 9.4, (9.4.3), p.~214. Original book scan inspected at \url{https://sites.math.rutgers.edu/~zeilberg/EM18/TitchmarshZeta.pdf}.
\bibitem{cc} A.~Connes and C.~Consani, \emph{Weil positivity and trace formula, the archimedean place}, supplied source \path{LIT_Connes_Consani_Weil_positivity_archimedean/weil-compo.tex}. Introduction, Theorem 1 (label \texttt{mainthmintro}); Appendix A for the Mellin normalization; Appendix C, Proposition C.1 (label \texttt{mainprop}), equation \texttt{appendweilnegav}. Section-based numbering must be retained.
\bibitem{li} X.-J.~Li, \emph{The positivity of a sequence of numbers and the Riemann hypothesis}, J. Number Theory \textbf{65} (1997), 325--333, \S1, Theorem 1. \url{https://www.sciencedirect.com/science/article/pii/S0022314X97921375}. Statement checked in a reproduction of the original article; no independent re-proof of the converse criterion is claimed.
\bibitem{baez} L.~B\'aez-Duarte, \emph{A strengthening of the Nyman--Beurling criterion for the Riemann hypothesis, 2}, arXiv:math/0205003v1, Theorem 1.1, p.~2. \url{https://arxiv.org/pdf/math/0205003}.
\bibitem{suzuki} M.~Suzuki, \emph{A canonical system of differential equations arising from the Riemann zeta-function}, arXiv:1204.1827v2, 23 September 2016. Proposition 1.2; (1.3), (1.7), (1.11); Theorem 2.3; \S5. \url{https://arxiv.org/pdf/1204.1827}.
\bibitem{conreyli} J.~B.~Conrey and X.-J.~Li, \emph{A note on some positivity conditions related to zeta and $L$-functions}, IMRN (2000), no.~18, 929--940. Theorems 1--2 and \S3.1, (3.1)--(3.4). \url{https://arxiv.org/pdf/math/9812166}.
\bibitem{connes} A.~Connes, \emph{Trace formula in noncommutative geometry and the zeros of the Riemann zeta function}, Selecta Math. (N.S.) \textbf{5} (1999), 29--106. Author text: \S V, Theorem 3; \S VII, Theorem 4; \S VIII, Theorem 5 and (16), (31)--(37), paragraph following (37); Appendix II, Theorem 6. Numbering here follows \url{https://alainconnes.org/wp-content/uploads/selecta.ps-2.pdf}.
\bibitem{deninger} C.~Deninger, \emph{Arithmetic geometry and analysis on foliated spaces}, 23 May 2005. \S3, (2), (3), (5)--(11), Proposition 3.2 and Theorem 3.3; cup/star calculation pp.~8--9. \url{https://swc-math.github.io/dls/DLSDeninger.pdf}.
\bibitem{sga5} A.~Grothendieck and collaborators, \emph{Cohomologie $\ell$-adique et fonctions $L$}, SGA 5, LNM 589. Expos\'e III, Corollary 4.7, (4.7.1), and Corollary 4.8, printed p.~109; Expos\'e III B, Corollary 1.3.1, printed p.~143. Original scans inspected: \url{https://library.slmath.org/nonmsri/sga/sga/5/SGA5-page-121.png} and \url{https://library.slmath.org/nonmsri/sga/sga/5/SGA5-page-155.png}.
\bibitem{sga6} P.~Berthelot, A.~Grothendieck, L.~Illusie and collaborators, \emph{Th\'eorie des intersections et th\'eor\`eme de Riemann--Roch}, SGA 6, LNM 225. Expos\'e V, Definition 2.1, Example 2.2(b), Lemma 2.2.1, (7.1.1); exact text inspected at \url{https://grothendiecksga.com/read/sga6/fr/V-2.html} and \url{https://grothendiecksga.com/read/sga6/fr/V-7.html}.
\bibitem{weil1} P.~Deligne, \emph{La conjecture de Weil. I}, Publ. Math. IHES \textbf{43} (1974), 273--307. (1.5.1), (1.5.3), (1.5.4), Theorem (1.6), pp.~275--276. \url{https://www.numdam.org/item/PMIHES_1974__43__273_0.pdf}.
\bibitem{ramare} O.~Ramar\'e, \emph{Corrigendum to ``Explicit estimates for the summatory function of $\Lambda(n)/n$ from the one of $\Lambda(n)$''}, corrected Lemma 2.1, p.~1, and its integral expression, p.~2. Corrects Lemma 2.2 of Acta Arith. \textbf{159} (2013), no.~2, 113--122. \url{https://ramare-olivier.github.io/Maths/Corrigendum-PsiTildeSimple-Acta-02.pdf}.
\bibitem{neretin} Yu.~A.~Neretin, \emph{Matrix balls, radial analysis of Berezin kernels, and hypergeometric determinants}, arXiv:math/0012220v2 (3 April 2001); Mosc. Math. J. \textbf{1} (2001), 157--220. Theorem 3.8, preprint p.~11; \S4.1 (4.1), p.~15; \S4.4 (4.7)--(4.9), p.~16. \url{https://arxiv.org/pdf/math/0012220v2}.
\bibitem{mcmullen} C.~McMullen, \emph{Advanced Complex Analysis}, Harvard Math 213a course notes, version dated 16 November 2026. Theorem 1.18, p.~19 (Liouville); Theorem 3.16, p.~92 (finite-order Hadamard factorization). \url{https://people.math.harvard.edu/~ctm/home/text/class/harvard/101/22/html/www/math213a/home/course/course.pdf}.
\end{thebibliography}
\end{document}
